% Results and discussion section for the FeRRy paper (draft, 8 Oct 2026); the
% setup and the design of each study are in setup.tex. Assembled into the
% full paper by DeveloperDocs/paper/assemble.py, which replaces each
%   %%TABLE:<label>%%  and  %%FIGURE:<label>%%
% line with the block of that label from results/exp5/paper/exp5_tables.tex or
% exp5_figures.tex (generated from the scores). Numbers in the prose come from
% results/exp5/scores (the FQ-vs-E3 paragraph from
% results/exp5/scores/b2/exploratory_fq_vs_e3.md) and the FerrySim reports.
%%%%% BEGIN SECTION %%%%%
\section{Evaluation Results and Discussion}
\label{sec:results}

\subsection{Against the State of the Art}
\label{sec:res:sota}

With one mule at the knee budget, F reaches $\tau$ in 194\,s and each of the
seven baselines needs 264--282\,s (Table~\ref{tab:exp5_headline}): F is
27--31\% faster, every comparison is a claim, and F gets there first in
80--90\% of paired seeds. FedEx with its own merge (D4$_{\text{FedEx}}$) is as
slow (281\,s) but misses $\tau$ once, so its comparison is not a claim. FX
(178\,s) and FQ (180\,s) do not differ from F; FQ, flown later, also beats D4
and E3 at the knee in its exploratory family. Under the stress budget only the
FedEx route is significantly slower: the other baselines trail by 28--60\,s but
stay within the budget by leaving devices out. With three mules every scheduler
reaches $\tau$ in 52--65\,s except D4$_{\text{FedEx}}$ (101\,s), whose
unweighted, one-tour-old updates slow learning. Without a mule (H0), FL over the
degraded link merges 0.5 updates per round and ends at 0.686 accuracy, barely
above the all-benign 0.667, against F's 3.5 and 0.831.

%%TABLE:tab:exp5_headline%%

\subsection{Why FeRRy Is Faster: Reach}
\label{sec:res:mechanism}

F does not need fewer rounds (every arm reaches $\tau$ in 1.65--1.84 rounds,
with near-identical update yield); its rounds are shorter
(Fig.~\ref{fig:exp5_mechanism}). Every baseline flies the wide band, which
reaches a device only from nearby, and spends about 140\,s of a mission in
transit. F's plan chooses the narrow band in 89\% of missions, whose range
covers the field from about one stop: F spends 15\,s in transit and 107\,s
dwelling, and its missions are 22--26\% shorter. FX and FQ inherit this reach
(narrow or medium on 95\% and 98\% of missions; Fig.~\ref{fig:exp5_learning}b).
Pinned to the wide band, F reaches $\tau$ in 265\,s, D4's own time; pinned to
narrow, in 196\,s. (These times are descriptive: the band-pinning study's
pre-registered metric, the age of updates, does not differ.) The choice is near
the best, serving as many devices as an exhaustive oracle at the knee budget and
4 points fewer under stress, and it adapts (Fig.~\ref{fig:exp5_mechanism}c): as the
field fills and the budget tightens, F shifts to medium and wide, mostly wide at
$N=24$ under stress.

%%FIGURE:fig:exp5_mechanism%%

\subsection{Scale and Decision Cost}
\label{sec:res:scale}

At $N=24$ (Table~\ref{tab:exp5_scale}), F is significantly faster than H1 and D3
at the knee (749 against 951 and 911\,s), ties the FedEx route (752\,s), and
beats it under the stress budget (661 against 752\,s); at $N=12$ no comparison
is a claim, and FQ (461 and 389\,s) does not differ from FX (424 and 386\,s).
Planning stays cheap (Table~\ref{tab:exp5_cost}): alone on one host, a plan
takes 0.04\,s at $N=6$ and 1.5\,s at $N=96$ (p95 2.0\,s). In FerrySim the plan
serves 54--57\% of the devices at each scale's knee budget up to $N=96$, and at
$N=12$ more mules cut F's time to $\tau$ from 486\,s to 226\,s with two and
151\,s with three (Table~\ref{tab:exp5_mules}).

%%TABLE:tab:exp5_scale%%

%%TABLE:tab:exp5_cost%%

%%TABLE:tab:exp5_mules%%

\subsection{The Design Claims}
\label{sec:res:claims}

Five tests hold (Table~\ref{tab:exp5_claims}). Without the coverage term the
round-close rate falls from 99\% to 74\% (C2). Under the stress budget FeRRy's
age of updates is 0.78 missions, against 1.06 for MAX-AoI, 1.07 for the Whittle
index and 1.10 without the coverage term (C5). On the time-varying backhaul the
adaptive controller lowers it from 0.88 to 0.69, close to the most adaptation
can give, since the controller sees the SNR that decides each loss. Other parts
show no measurable effect here: pinning the band leaves the age of updates
unchanged (its effect is on time), the dwell term and the age cap change nothing
over four missions, and the merge rules tie because collected updates almost
always have merge age zero (FedBuff's buffering slows learning, and FedProx is
not significantly faster). The selection literature's deadline forms miss
slightly more deadlines (47\% against 42\% under stress), but not significantly.

%%TABLE:tab:exp5_claims%%

\subsection{Learning Inside the Plan}
\label{sec:res:rl}

The learned-score study is flat: no $\gamma$ beats $\gamma=0$ by the
pre-registered margin, and the best ($\gamma=0.75$) matches FX while trailing a
one-step greedy rule. In the full system FQ is again indistinguishable from FX
(C4) and, at $N=6$, from F, so FX is FeRRy's per-stop rule. The two decide on
different information (FQ before serving, from the priced dwell; FX after, from
the realized mission time), so this compares two rules, not two rankings.

The exploratory comparison sets FQ against E3, a single-agent DQN after Chen et
al.~\cite{chen2023model} that picks every next stop with no plan, on the same 20
seeds (Fig.~\ref{fig:exp5_learning}). At the knee FQ reaches $\tau$ in 180\,s
against E3's 280\,s (difference 100\,s, 95\% CI [63, 141], Holm $p<0.001$ across
the eight exploratory comparisons); under stress it is faster over 19 shared
seeds (174 against 211\,s), but the CI includes zero. FQ inherits the plan's
band and spends 15\,s of a 155\,s mission in transit, against E3's 144\,s of
228\,s, and in FerrySim E3's held-out return falls 6\% below FX's while FQ's
matches it. The gain comes from the plan, not from learning, which supports the
two-clock structure (C4).

%%FIGURE:fig:exp5_learning%%

\subsection{Robustness and Limitations}
\label{sec:res:robust}

With strongly non-IID data (Dirichlet $\alpha=0.1$;
Table~\ref{tab:exp5_robust}), every arm reaches $\tau$ in about half the trials,
and F stays significantly faster than H1 and D3 where both do. The radio study
flies the time-varying backhaul, so its default (F at 216\,s) is not the main
configuration (194\,s); under its harsher channels F stays ahead of H1 without a
claim, and FX is significantly faster than F on the default channel and with
8\,dB shadowing. Two limitations stand out. FeRRy's plan prices the mean SNR, so
under the stress budget 30\% of its missions overrun, by 28\,s on average
(median 20\,s; Table~\ref{tab:exp5_headline}); the FedEx route, which never skips
a device, overruns in 85\%, by 22\,s on average, and F keeps nearly its update
yield with fewer stops. And the
plan does not price training time: when devices take time to train, F's yield
falls from 3.5 to 0.7 updates per round and it reaches $\tau$ in only 70--75\% of
trials, while H1, D4 and D5 keep 85--100\%.

%%TABLE:tab:exp5_robust%%
%%%%% END SECTION %%%%%
